\documentclass[a4paper,12pt]{article}

\usepackage{amsmath,amssymb}
\usepackage{graphicx}
\usepackage{setspace}
\usepackage{float}
\usepackage{enumitem}
\usepackage{hyperref}
\usepackage{xurl}
\usepackage{geometry}
\usepackage{fancyhdr}
\usepackage{array}
\usepackage{tabularx}
\usepackage{multirow}
\usepackage{booktabs}
\usepackage{xcolor}
\usepackage{longtable}
\usepackage{tikz}
\usetikzlibrary{shapes,arrows.meta,positioning,calc,fit,decorations.pathreplacing,backgrounds}

\geometry{top=2.5cm,bottom=2.5cm,left=1.8cm,right=1.8cm}
\setstretch{1.55}
\setlength{\parskip}{0.35em}
\setlength{\parindent}{0.8cm}
\setlength{\headheight}{22pt}
\Urlmuskip=0mu plus 1mu

\definecolor{myblue}{RGB}{70,130,180}
\definecolor{mygreen}{RGB}{60,179,113}
\definecolor{myred}{RGB}{220,20,60}
\definecolor{myorange}{RGB}{255,140,0}
\definecolor{mypurple}{RGB}{147,112,219}

\usepackage{xepersian}
\IfFontExistsTF{XB Niloofar}{
  \settextfont{XB Niloofar}
  \setdigitfont{XB Niloofar}
}{
  \IfFontExistsTF{Amiri}{
    \settextfont{Amiri}
    \setdigitfont{Amiri}
  }{
    \settextfont{Noto Naskh Arabic}
    \setdigitfont{Noto Naskh Arabic}
  }
}
\IfFontExistsTF{Times New Roman}{\setlatintextfont{Times New Roman}}{\setlatintextfont{DejaVu Serif}}

\hypersetup{
  unicode=true,
  breaklinks=true,
  colorlinks=true,
  linkcolor=blue!45!black,
  urlcolor=blue!45!black
}

\newcommand{\EN}[1]{\lr{#1}}
\newcommand{\NPone}{\lr{NoiseProj1}}
\newcommand{\NPtwo}{\lr{NoiseProj2}}
\newcommand{\InfoNCE}{\lr{InfoNCE}}
\newcommand{\QR}{\lr{QR}}
\newcommand{\TopK}{\lr{Top-k}}
\newcommand{\MSE}{\operatorname{MSE}}
\newcommand{\EMA}{\lr{EMA}}
\newcommand{\Dterm}{\lr{D}}
\newcommand{\Iterm}{\lr{I}}
\newcommand{\Nterm}{\lr{N}}
\newcommand{\Sterm}{\lr{S}}

\newcolumntype{C}[1]{>{\centering\arraybackslash}m{#1}}
\newcolumntype{R}[1]{>{\raggedleft\arraybackslash}m{#1}}
\newcolumntype{L}[1]{>{\raggedright\arraybackslash}m{#1}}

\pagestyle{fancy}
\fancyhf{}
\fancyhead[R]{گزارش فنی روش‌های فعلی کاهش‌بُعد آگاه از نویز}
\fancyhead[L]{\thepage}
\renewcommand{\headrulewidth}{0.4pt}

\begin{document}

\begin{titlepage}
\thispagestyle{empty}
\centering
\vspace*{2.7cm}
{\Large\bfseries باسمه تعالی\par}
\vspace{1.8cm}
{\Huge\bfseries روش‌های فعلی کاهش‌بُعد آگاه از نویز\par}
\vspace{0.5cm}
{\Large\bfseries توضیح کامل \NPone{}، \NPtwo{} و خانواده‌های ترکیبی فعال\par}
\vspace{2cm}
\begin{minipage}{0.88\textwidth}
\centering
این متن فقط تعریف فنی روش‌های فعال پروژه را ارائه می‌کند. روش‌های کنارگذاشته‌شده، گزارش نتایج آزمایش‌ها و روند تاریخی توسعه در آن وارد نشده‌اند. تمرکز متن بر آخرین تعریف ریاضی و پیاده‌سازی صحیح دو روش اصلی و خانواده‌های فعلی آن‌ها است.
\end{minipage}
\vfill
{\large گزارش فنی مستقل\par}
\end{titlepage}

\tableofcontents
\newpage

\section{صورت مسئله و قرارداد نمادگذاری}

هدف، یادگیری یک نگاشت خطی از فضای تعبیه‌ی اولیه‌ی $d$-بعدی به فضای کم‌بُعد $r$-بعدی است، به‌گونه‌ای که اطلاعات معنایی موردنیاز برای بازیابی تا حد امکان حفظ شود و حساسیت بازنمایی به نویزهای متنی کاهش یابد. مدل تعبیه‌ساز ثابت است و تنها ماتریس کاهش‌بُعد آموخته یا محاسبه می‌شود.

برای متن تمیز $t_i$ و نمای نویزی $v$-ام آن، داریم:
\begin{equation}
 x_i=f(t_i)\in\mathbb{R}^{d},
 \qquad
 \tilde{x}_i^{(v)}=f\!\left(\tilde{t}_i^{(v)}\right)\in\mathbb{R}^{d},
\end{equation}
که در آن $f$ مدل تعبیه‌ساز ثابت است. در کد، بردارها سطری نگهداری می‌شوند و نگاشت کم‌بُعد به‌صورت زیر اجرا می‌شود:
\begin{equation}
 z=xW,
 \qquad
 W\in\mathbb{R}^{d\times r},
 \qquad
 r<d.
\end{equation}

\begin{center}
\begin{tabularx}{0.95\textwidth}{|C{2.4cm}|X|}
\hline
\textbf{نماد} & \textbf{تعریف} \\
\hline
$X\in\mathbb{R}^{N\times d}$ & ماتریس تعبیه‌های تمیز اسناد آموزشی \\
\hline
$\widetilde{X}^{(v)}$ & ماتریس تعبیه‌های نمای نویزی $v$-ام \\
\hline
$E^{(v)}=\widetilde{X}^{(v)}-X$ & جابه‌جایی ایجادشده بر اثر نویز \\
\hline
$W\in\mathbb{R}^{d\times r}$ & ماتریس کاهش‌بُعد \\
\hline
$\Sigma_s$ & کوواریانس سیگنال، برآوردشده از تعبیه‌های تمیز \\
\hline
$\Sigma_n$ & کوواریانس نویز، برآوردشده از اختلاف تعبیه‌های نویزی و تمیز \\
\hline
$B=\Sigma_n+\rho I$ & متریک نویز منظم‌شده \\
\hline
\Dterm{} & ترم نویززدایی \\
\hline
\Iterm{} & ترم ناوردایی تقابلی \\
\hline
\Nterm{} & ترم سرکوب نویز \\
\hline
\Sterm{} & ترم حفظ سیگنال \\
\hline
\end{tabularx}
\end{center}

در آموزش، متن‌های پیکره برای ساخت نماهای تمیز و نویزی استفاده می‌شوند. کوئری‌های ارزیابی، نسخه‌های نویزی آن‌ها و برچسب‌های ارتباط وارد یادگیری ماتریس کاهش‌بُعد نمی‌شوند.

\section{روش ۱: یادگیری گرادیانی چندهدفه‌ی \NPone{}}

\subsection{هدف روش}

روش ۱ یک نگاشت خطی قابل‌آموزش است که به‌طور هم‌زمان پنج هدف را دنبال می‌کند:

\begin{enumerate}[label=\arabic*)]
 \item بازسازی تعبیه‌ی تمیز از روی تعبیه‌ی نویزی؛
 \item نزدیک‌کردن نماهای نویزی متن یکسان در فضای کم‌بُعد؛
 \item کاهش انرژی جهت‌های حساس به نویز؛
 \item جلوگیری از حذف بیش از حد تغییرات مفید موجود در داده‌ی تمیز؛
 \item کنترل هندسه‌ی ماتریس نگاشت با متعامدسازی اقلیدسی.
\end{enumerate}

روش ۱ با بهینه‌سازی گرادیانی و به‌کمک الگوریتم \EN{Adam} آموزش داده می‌شود. مدل تعبیه‌ساز در این فرایند ثابت می‌ماند.

\subsection{ساخت نماهای آموزشی}

برای هر متن آموزشی، چند نمای نویزی مستقل ساخته می‌شود. تنظیم پیش‌فرض روش شامل چهار نما با شدت نویز $0.15$ و خانواده‌های نویز تایپی و خطای مشابه \EN{OCR} است. در هر نما، خانواده‌ی نویز برای هر متن به‌صورت مستقل انتخاب می‌شود. بنابراین یک نمای کامل لزوماً فقط از یک نوع نویز تشکیل نشده و می‌تواند نمونه‌هایی از چند خانواده‌ی نویز داشته باشد.

مجموعه‌ی آموزشی به‌صورت زیر است:
\begin{equation}
 \mathcal{V}_i=
 \left\{
 x_i,
 \tilde{x}_i^{(1)},
 \ldots,
 \tilde{x}_i^{(V)}
 \right\}.
\end{equation}

در هر بار دریافت نمونه از داده، دو نمای نویزی متفاوت از همان متن انتخاب می‌شوند:
\begin{equation}
 \left(\tilde{x}_i^{(a)},\tilde{x}_i^{(b)},x_i\right),
 \qquad a\neq b.
\end{equation}
انتخاب جفت برای هر شاخص داده قطعی و وابسته به بذر تصادفی است؛ در نتیجه، زمان‌بندی کارگرهای \EN{DataLoader} نمی‌تواند جفت‌های مثبت را تغییر دهد. در عین حال، جابه‌جایی ترتیب نمونه‌ها در هر دوره‌ی آموزش حفظ می‌شود.

\subsection{نگاشت و بازسازی در روش ۱}

ماتریس $W$ در روش ۱ با مقداردهی اولیه‌ی متعامد ساخته می‌شود:
\begin{equation}
 W_0^{T}W_0=I_r.
\end{equation}

نمایش کم‌بُعد برابر است با:
\begin{equation}
 z=xW.
\end{equation}

بازسازی روش ۱ از وزن‌های بسته‌شده استفاده می‌کند:
\begin{equation}
 \hat{x}=zW^{T}=xWW^{T}.
\end{equation}

رابطه‌ی $WW^{T}$ زمانی پروژکتور متعامد دقیق است که ستون‌های $W$ متعامد باشند. روش ۱ این شرط را هم با جریمه‌ی نرم و هم با بازگردانی سخت دوره‌ای حفظ می‌کند.

\subsection{ترم نویززدایی \Dterm{}}

برای دو نمای نویزی انتخاب‌شده، بازسازی‌ها عبارت‌اند از:
\begin{equation}
 \hat{x}_i^{(a)}=\tilde{x}_i^{(a)}WW^{T},
 \qquad
 \hat{x}_i^{(b)}=\tilde{x}_i^{(b)}WW^{T}.
\end{equation}

زیان نویززدایی، میانگین دو خطای بازسازی نسبت به تعبیه‌ی تمیز است:
\begin{equation}
 \mathcal{L}_{\mathrm{denoise}}
 =
 \frac{1}{2}
 \left[
 \MSE\!\left(\hat{x}^{(a)},x\right)
 +
 \MSE\!\left(\hat{x}^{(b)},x\right)
 \right].
\end{equation}

این ترم روش را وادار می‌کند اطلاعات مشترک میان متن تمیز و نسخه‌های نویزی را در زیرفضای کم‌بُعد نگه دارد. هدف بازسازی، بردار نویزی نیست؛ بلکه بردار تمیز متناظر است.

\subsection{ترم ناوردایی تقابلی \Iterm{}}

دو نمای نویزی به فضای کم‌بُعد نگاشت می‌شوند:
\begin{equation}
 z_i^{(a)}=\tilde{x}_i^{(a)}W,
 \qquad
 z_i^{(b)}=\tilde{x}_i^{(b)}W.
\end{equation}

پس از نرمال‌سازی $\ell_2$، ماتریس شباهت یک دسته برابر است با:
\begin{equation}
 S_{ij}
 =
 \frac{
 \left\langle
 \bar{z}_i^{(a)},
 \bar{z}_j^{(b)}
 \right\rangle
 }{\tau},
\end{equation}
که $\tau$ دمای زیان تقابلی است. نمونه‌ی مثبت برای سطر $i$، نمای دیگر همان متن یعنی ستون $i$ است و سایر نمونه‌های دسته نقش منفی را دارند. زیان متقارن \InfoNCE{} به‌صورت زیر محاسبه می‌شود:
\begin{equation}
 \mathcal{L}_{\mathrm{inv}}
 =
 \frac{1}{2}
 \left[
 \mathrm{CE}(S,\mathbf{y})
 +
 \mathrm{CE}(S^{T},\mathbf{y})
 \right],
 \qquad
 \mathbf{y}=(0,1,\ldots,m-1).
\end{equation}

حالت پیش‌فرض، مقایسه‌ی نویزی--نویزی است. پیاده‌سازی امکان مقایسه‌ی تمیز--نویزی یا ترکیب هر دو نوع را نیز دارد، ولی تعریف استاندارد روش ۱ از حالت نویزی--نویزی استفاده می‌کند.

\subsection{برآورد برخط کوواریانس نویز}

برای هر دسته، دو مجموعه‌ی اختلاف نویز ساخته می‌شوند:
\begin{equation}
 E^{(a)}=\widetilde{X}^{(a)}-X,
 \qquad
 E^{(b)}=\widetilde{X}^{(b)}-X.
\end{equation}

پس از اتصال دو مجموعه و مرکزکردن آن‌ها، کوواریانس دسته به‌صورت زیر برآورد می‌شود:
\begin{equation}
 \widehat{\Sigma}_{n}^{\mathrm{batch}}
 =
 \frac{1}{M-1}E_c^{T}E_c.
\end{equation}

کوواریانس سراسری با میانگین متحرک نمایی به‌روزرسانی می‌شود:
\begin{equation}
 \Sigma_n
 \leftarrow
 \mu\Sigma_n
 +(1-\mu)\widehat{\Sigma}_{n}^{\mathrm{batch}}.
\end{equation}
در تنظیم پیش‌فرض $\mu=0.99$ است. به‌روزرسانی فقط در فازهایی انجام می‌شود که وزن ترم سرکوب نویز غیرصفر باشد.

\subsection{ترم سرکوب نویز \Nterm{}}

انرژی نویز پس از نگاشت با عبارت زیر اندازه‌گیری می‌شود:
\begin{equation}
 \mathcal{L}_{\mathrm{suppress}}
 =
 \frac{1}{r}
 \operatorname{Tr}\!\left(W^{T}\Sigma_nW\right).
\end{equation}

کمینه‌کردن این ترم، ستون‌های $W$ را از جهت‌هایی دور می‌کند که تغییرات ناشی از نویز در آن‌ها زیاد است.

\subsection{ترم حفظ سیگنال \Sterm{}}

نمایش تمیز کم‌بُعد برابر است با:
\begin{equation}
 Z=XW.
\end{equation}
پس از مرکزکردن ستون‌های $Z$، میانگین واریانس حفظ‌شده محاسبه می‌شود. چون تابع هدف کمینه می‌شود، ترم حفظ سیگنال با علامت منفی وارد می‌شود:
\begin{equation}
 \mathcal{L}_{\mathrm{signal}}
 =
 -\frac{1}{Nr}
 \left\|Z-\mathbf{1}\bar{z}^{T}\right\|_{F}^{2}.
\end{equation}

این ترم از فروپاشی زیرفضا و حذف کامل تغییرات مفید داده‌ی تمیز جلوگیری می‌کند. در روش ۱، متعامدبودن و نرمال‌بودن ستون‌های $W$ مانع رشد نامحدود این عبارت می‌شود.

\subsection{ترم متعامدسازی اقلیدسی}

جریمه‌ی نرم هندسی روش ۱ برابر است با:
\begin{equation}
 \mathcal{L}_{\mathrm{ortho}}
 =
 \left\|W^{T}W-I_r\right\|_{F}^{2}.
\end{equation}

این جریمه انحراف تدریجی ستون‌های $W$ از پایه‌ی متعامد را کاهش می‌دهد. علاوه بر آن، پس از هر تعداد مشخصی گام بهینه‌سازی، یک \QR{} سخت اجرا می‌شود:
\begin{equation}
 W=QR,
 \qquad
 W\leftarrow Q.
\end{equation}
در تنظیم پیش‌فرض، این بازگردانی هر $25$ گام انجام می‌شود.

\subsection{تابع هدف کامل روش ۱}

تابع هدف استاندارد روش ۱ عبارت است از:
\begin{equation}
 \mathcal{L}_{\mathrm{NP1}}
 =
 \lambda_D\mathcal{L}_{\mathrm{denoise}}
 +
 \lambda_I\mathcal{L}_{\mathrm{inv}}
 +
 \lambda_N\mathcal{L}_{\mathrm{suppress}}
 +
 \lambda_S\mathcal{L}_{\mathrm{signal}}
 +
 \lambda_O\mathcal{L}_{\mathrm{ortho}}.
\end{equation}

مقادیر پیش‌فرض پیاده‌سازی فعلی:
\begin{center}
\begin{tabular}{|C{4.1cm}|C{3cm}|C{4.5cm}|}
\hline
\textbf{پارامتر} & \textbf{مقدار} & \textbf{نقش} \\
\hline
$\lambda_D$ & $1.0$ & بازسازی تعبیه‌ی تمیز \\
\hline
$\lambda_I$ & $0.5$ & ناوردایی تقابلی \\
\hline
$\lambda_N$ & $0.3$ & سرکوب نویز \\
\hline
$\lambda_S$ & $0.1$ & حفظ سیگنال \\
\hline
$\lambda_O$ & $10^{-3}$ & متعامدسازی نرم \\
\hline
$\tau$ & $0.1$ & دمای \InfoNCE{} \\
\hline
نرخ یادگیری & $5\times10^{-4}$ & بهینه‌ساز \EN{Adam} \\
\hline
\end{tabular}
\end{center}

\subsection{زمان‌بندی فازهای آموزش}

روش ۱ همه‌ی ترم‌ها را از نخستین دوره با وزن کامل فعال نمی‌کند. زمان‌بندی پیش‌فرض سه مرحله دارد:

\begin{enumerate}[label=\arabic*)]
 \item \textbf{گرم‌سازی، ۲۰ درصد نخست:}
 \begin{equation}
 \lambda_I=0,
 \qquad
 \lambda_N=0,
 \end{equation}
 در حالی که نویززدایی، حفظ سیگنال و متعامدسازی فعال‌اند.

 \item \textbf{مرحله‌ی ناوردایی، از ۲۰ تا ۵۰ درصد آموزش:}
 وزن ناوردایی و حفظ سیگنال با ضریب $0.5$ اعمال می‌شود و سرکوب نویز هنوز غیرفعال است.

 \item \textbf{مرحله‌ی کامل، نیمه‌ی دوم آموزش:}
 همه‌ی ترم‌ها با وزن اصلی خود فعال می‌شوند.
\end{enumerate}

پس از هر گام، محدودیت سخت هندسی در بازه‌های تعیین‌شده اعمال می‌شود. اگر زیان نامتناهی یا غیرعددی شود، آموزش با خطای صریح متوقف می‌شود.

\subsection{پیاده‌سازی نهایی و اصلاحات فنی روش ۱}

پیاده‌سازی فعلی روش ۱ دارای اصلاحات زیر است:

\begin{itemize}
 \item قرارداد بردارهای سطری و نگاشت $xW$ در همه‌ی ماژول‌ها یکسان شده است.
 \item شکل و متناهی‌بودن همه‌ی ماتریس‌های ورودی، وزن اولیه و کوواریانس‌ها کنترل می‌شود.
 \item انتخاب دو نمای مثبت برای هر نمونه قطعی است و به زمان‌بندی پردازش موازی وابسته نیست.
 \item سازنده‌ی نویز، اشیای نویزساز را در هر نما دوباره استفاده می‌کند و از ساخت تکراری و پرهزینه‌ی آن‌ها برای هر متن جلوگیری می‌شود؛ معنای آماری تولید نماها تغییر نمی‌کند.
 \item نماهای نویزی یکسان در طول یک اجرای آزمایش با \EN{TrainingViewStore} در حافظه بازاستفاده می‌شوند تا روش‌های مختلف دقیقاً روی داده‌ی آموزشی یکسان اجرا شوند.
 \item کنترل‌های هندسی به‌صورت ماژولار پیاده‌سازی شده‌اند، ولی تنظیم پیش‌فرض روش ۱ همچنان بازسازی بسته‌شده، جریمه‌ی متعامدسازی اقلیدسی و \QR{} دوره‌ای است.
 \item گزینه‌های برش گرادیان و کاهش وزن در کد وجود دارند، اما در تعریف پیش‌فرض روش ۱ غیرفعال‌اند.
\end{itemize}

\section{روش ۲: حل بسته‌ی نسبت سیگنال به نویز \NPtwo{}}

\subsection{هدف روش}

روش ۲ به‌جای آموزش گرادیانی چندترمی، زیرفضای کم‌بُعد را مستقیماً از آمار مرتبه‌ی دوم داده‌های تمیز و اغتشاش‌های نویزی استخراج می‌کند. هدف انتخاب جهت‌هایی است که در آن‌ها واریانس سیگنال زیاد و واریانس نویز کم باشد.

این روش پارامترهای قابل‌آموزش تکرارشونده ندارد و پس از محاسبه‌ی کوواریانس‌ها، با حل یک مسئله‌ی مقدارویژه‌ی تعمیم‌یافته به جواب می‌رسد.

\subsection{برآورد کوواریانس سیگنال}

میانگین تعبیه‌های تمیز برابر است با:
\begin{equation}
 \mu_s=\frac{1}{N}\sum_{i=1}^{N}x_i.
\end{equation}

کوواریانس سیگنال از داده‌های مرکز‌شده محاسبه می‌شود:
\begin{equation}
 \Sigma_s
 =
 \frac{1}{N-1}
 \sum_{i=1}^{N}
 (x_i-\mu_s)(x_i-\mu_s)^{T}.
\end{equation}

مرکزکردن سیگنال در تنظیم پیش‌فرض فعال است.

\subsection{برآورد کوواریانس نویز}

برای هر نمای نویزی، اختلاف با بردار تمیز متناظر ساخته می‌شود:
\begin{equation}
 e_i^{(v)}=\tilde{x}_i^{(v)}-x_i.
\end{equation}

تمام اختلاف‌ها از نماهای مختلف روی هم قرار می‌گیرند. اگر $V$ نمای نویزی وجود داشته باشد، مجموعه‌ی نویز شامل $NV$ بردار است. میانگین نویز و کوواریانس آن به‌صورت زیر هستند:
\begin{equation}
 \mu_n
 =
 \frac{1}{NV}
 \sum_{i=1}^{N}\sum_{v=1}^{V}e_i^{(v)},
\end{equation}
\begin{equation}
 \Sigma_n
 =
 \frac{1}{NV-1}
 \sum_{i=1}^{N}\sum_{v=1}^{V}
 \left(e_i^{(v)}-\mu_n\right)
 \left(e_i^{(v)}-\mu_n\right)^{T}.
\end{equation}

مرکزکردن نویز نیز در تنظیم پیش‌فرض فعال است.

\subsection{انقباض کوواریانس‌ها}

برای بهبود پایداری عددی، هر کوواریانس می‌تواند به سمت ماتریس همانی مقیاس‌شده منقبض شود:
\begin{equation}
 \Sigma'
 =
 (1-\alpha)\Sigma
 +
 \alpha\frac{\operatorname{Tr}(\Sigma)}{d}I.
\end{equation}

در تنظیم پیش‌فرض:
\begin{equation}
 \alpha_s=0,
 \qquad
 \alpha_n=0.05.
\end{equation}
بنابراین کوواریانس سیگنال بدون انقباض و کوواریانس نویز با انقباض ملایم استفاده می‌شود.

\subsection{متریک نویز منظم‌شده}

برای جلوگیری از تکین‌شدن یا بدشرط‌شدن کوواریانس نویز، ماتریس زیر ساخته می‌شود:
\begin{equation}
 B=\Sigma_n+\rho I.
\end{equation}

مقدار $\rho$ نسبت به مقیاس متوسط واریانس نویز تعیین می‌شود:
\begin{equation}
 \rho
 =
 \max\left(
 \varepsilon,
 \gamma\max\left(\frac{\operatorname{Tr}(\Sigma_n)}{d},\varepsilon\right)
 \right),
\end{equation}
که در آن $\gamma$ پارامتر منظم‌سازی و $\varepsilon$ کران عددی کوچک است. مقادیر پیش‌فرض عبارت‌اند از:
\begin{equation}
 \gamma=10^{-3},
 \qquad
 \varepsilon=10^{-8}.
\end{equation}

\subsection{مسئله‌ی بهینه‌سازی}

برای یک جهت $w$، نسبت سیگنال به نویز برابر است با:
\begin{equation}
 R(w)
 =
 \frac{w^{T}\Sigma_sw}{w^{T}Bw}.
\end{equation}

بیشینه‌سازی این نسبت به مسئله‌ی مقدارویژه‌ی تعمیم‌یافته منجر می‌شود:
\begin{equation}
 \Sigma_sw
 =
 \lambda Bw.
\end{equation}

برای ساخت فضای $r$-بعدی، بردارهای متناظر با $r$ مقدارویژه‌ی بزرگ‌تر انتخاب می‌شوند. شکل ماتریسی هدف را می‌توان به‌صورت زیر نوشت:
\begin{equation}
 \max_{W\in\mathbb{R}^{d\times r}}
 \operatorname{Tr}\!\left(W^{T}\Sigma_sW\right),
 \qquad
 \text{با قید}
 \qquad
 W^{T}BW=I_r.
\end{equation}

\subsection{حل پایدار با سفیدسازی متریک نویز}

به‌جای حل مستقیم مسئله‌ی تعمیم‌یافته، ابتدا تجزیه‌ی ویژه‌ی $B$ محاسبه می‌شود:
\begin{equation}
 B=U\Lambda U^{T}.
\end{equation}

سپس معکوس ریشه‌ی دوم آن ساخته می‌شود:
\begin{equation}
 B^{-1/2}
 =
 U\Lambda^{-1/2}U^{T}.
\end{equation}
مقادیر ویژه از پایین با $\varepsilon$ کران‌گذاری می‌شوند تا تقسیم بر عدد بسیار کوچک رخ ندهد.

کوواریانس سیگنال در فضای سفیدشده برابر است با:
\begin{equation}
 C
 =
 B^{-1/2}\Sigma_sB^{-1/2}.
\end{equation}

اگر $V_r$ شامل $r$ بردار ویژه‌ی اصلی $C$ باشد، پایه‌ی اولیه‌ی روش ۲ برابر است با:
\begin{equation}
 W=B^{-1/2}V_r.
\end{equation}

\subsection{متعامدسازی در متریک نویز}

ستون‌های خروجی روش ۲ قرار نیست در هندسه‌ی اقلیدسی متعامد باشند. شرط درست آن‌ها این است:
\begin{equation}
 W^{T}BW=I_r.
\end{equation}

برای ترمیم خطای محاسبات ممیز شناور، ماتریس زیر ساخته می‌شود:
\begin{equation}
 G=W^{T}BW.
\end{equation}
سپس خروجی به‌صورت زیر اصلاح می‌شود:
\begin{equation}
 W\leftarrow WG^{-1/2}.
\end{equation}

این تبدیل گستره‌ی ستونی $W$ را عوض نمی‌کند و فقط شرط متعامدسازی متریکی را بازمی‌گرداند.

\subsection{چرا \QR{} اقلیدسی جزو روش ۲ نیست؟}

اجرای \QR{} معمولی روی $W$ شرط زیر را تحمیل می‌کند:
\begin{equation}
 W^{T}W=I_r.
\end{equation}
این شرط با قید اصلی روش ۲ متفاوت است. جایگزین‌کردن پایه‌ی متریکی با پایه‌ی اقلیدسی، مقیاس و هندسه‌ی مختصات کم‌بُعد را تغییر می‌دهد و می‌تواند شباهت کسینوسی پس از نگاشت را عوض کند. بنابراین، روش ۲ در پیاده‌سازی نهایی فقط در متریک $B$ بازمتعامد می‌شود و \QR{} اقلیدسی روی خروجی مستقل آن اجرا نمی‌شود.

\subsection{تبدیل نهایی}

پس از محاسبه‌ی $W$، هر بردار با یک ضرب ماتریسی ساده تبدیل می‌شود:
\begin{equation}
 z=xW.
\end{equation}

روش ۲ دوره‌ی آموزش، بهینه‌ساز، نرخ یادگیری و پس‌انتشار ندارد. تاریخچه‌ی آن شامل خلاصه‌ی مرحله‌ی حل بسته، مقادیر ویژه، نسبت هدف، خطای متعامدسازی متریکی و آمار کوواریانس‌ها است.

\subsection{پیاده‌سازی نهایی و اصلاحات فنی روش ۲}

اصلاحات و ویژگی‌های پیاده‌سازی فعلی عبارت‌اند از:

\begin{itemize}
 \item کوواریانس نویز از اختلاف جفت‌شده‌ی تعبیه‌ی نویزی و تمیز محاسبه می‌شود، نه از خود بردارهای نویزی.
 \item ماتریس‌های کوواریانس پیش از تجزیه به‌طور صریح متقارن می‌شوند.
 \item انقباض کوواریانس نویز و منظم‌سازی متناسب با مقیاس واریانس برای جلوگیری از بدشرطی اعمال می‌شود.
 \item مسئله‌ی تعمیم‌یافته با سفیدسازی متقارن حل می‌شود تا از وارون‌سازی مستقیم ماتریس جلوگیری شود.
 \item مقادیر ویژه‌ی کوچک در محاسبه‌ی ریشه‌ی معکوس کران‌گذاری می‌شوند.
 \item بازمتعامدسازی نهایی در متریک $B$ انجام می‌شود و گستره‌ی ستونی را حفظ می‌کند.
 \item نرمال‌سازی اقلیدسی ستون‌ها در تنظیم پیش‌فرض غیرفعال است، زیرا هندسه‌ی خروجی را تغییر می‌دهد.
 \item نماهای نویزی و تعبیه‌های تمیز می‌توانند با سایر روش‌ها در همان اجرای آزمایش بازاستفاده شوند.
\end{itemize}

\section{ابزارهای هندسی مشترک در خانواده‌های ترکیبی}

خانواده‌های فعلی با ترکیب مقداردهی اولیه، نوع بازسازی و قید هندسی ساخته می‌شوند. چهار ابزار اصلی آن‌ها در ادامه تعریف شده‌اند.

\subsection{بازسازی بسته‌شده}

وقتی ستون‌های $W$ اقلیدسی‌متعامدند، بازسازی زیر معتبر است:
\begin{equation}
 \hat{x}=xWW^{T}.
\end{equation}
این بازسازی در خانواده‌هایی استفاده می‌شود که \QR{} اقلیدسی سخت دارند.

\subsection{پروژکتور عمومی برای پایه‌ی غیرمتعامد}

اگر $W$ پایه‌ای تمام‌رتبه ولی غیرمتعامد باشد، عبارت $WW^{T}$ پروژکتور درست نیست. پروژکتور روی گستره‌ی ستونی $W$ برابر است با:
\begin{equation}
 P_W
 =
 W\left(W^{T}W+\epsilon I\right)^{-1}W^{T}.
\end{equation}

بازسازی به‌شکل زیر انجام می‌شود:
\begin{equation}
 \hat{x}=xP_W.
\end{equation}

در کد، معکوس ماتریس به‌طور صریح محاسبه نمی‌شود. ابتدا
\begin{equation}
 A=W^{T}W+\epsilon I
\end{equation}
ساخته می‌شود و ضرایب با \EN{torch.linalg.solve} به‌دست می‌آیند. این روش از نظر عددی پایدارتر از محاسبه‌ی مستقیم $A^{-1}$ است.

\subsection{بازگردانی اقلیدسی}

بازگردانی اقلیدسی با \QR{} انجام می‌شود:
\begin{equation}
 W=QR,
 \qquad
 W\leftarrow Q.
\end{equation}
پس از آن:
\begin{equation}
 W^{T}W=I_r.
\end{equation}

\subsection{بازگردانی متریکی}

برای حفظ هندسه‌ی روش ۲، ابتدا
\begin{equation}
 G=W^{T}BW
\end{equation}
محاسبه می‌شود و سپس:
\begin{equation}
 W\leftarrow WG^{-1/2}.
\end{equation}
این تبدیل شرط زیر را بازمی‌گرداند:
\begin{equation}
 W^{T}BW=I_r.
\end{equation}

در کد، $G$ متقارن می‌شود، تجزیه‌ی ویژه روی آن انجام می‌گیرد و مقادیر ویژه از پایین کران‌گذاری می‌شوند.

\subsection{کنترل مقیاس ستون‌ها}

وقتی \QR{} اقلیدسی حذف می‌شود، ترم حفظ سیگنال می‌تواند با بزرگ‌کردن نرم ستون‌ها مقدار منفی‌تری بگیرد. برای جلوگیری از رشد کنترل‌نشده، نرم هر ستون با نرم اولیه‌ی همان ستون مقایسه می‌شود:
\begin{equation}
 r_j
 =
 \frac{\lVert w_j\rVert_2}{\lVert w_{0,j}\rVert_2}.
\end{equation}

زیان کنترل مقیاس برابر است با:
\begin{equation}
 \mathcal{L}_{\mathrm{scale}}
 =
 \frac{1}{r}
 \sum_{j=1}^{r}
 \left(r_j^2-1\right)^2.
\end{equation}

رشد مرتبه‌ی چهار این زیان عمدی است، زیرا ترم منفی واریانس سیگنال با مقیاس به‌صورت درجه‌ی دوم رشد می‌کند.

\section{خانواده‌ی روش‌های ۴ تا ۷: روش ۱ با \QR{} اقلیدسی}

\subsection{تعریف خانواده}

این خانواده نسخه‌های مستقیم روش ۱ را بررسی می‌کند، بدون اینکه از روش ۲ برای مقداردهی اولیه استفاده شود. مسیر کلی آن عبارت است از:
\begin{equation}
 W_{\mathrm{random\;orthogonal}}
 \longrightarrow
 \text{آموزش گرادیانی روش ۱}
 \longrightarrow
 W_{\mathrm{final}}.
\end{equation}

ویژگی‌های مشترک خانواده:
\begin{itemize}
 \item مقداردهی اولیه‌ی تصادفی متعامد؛
 \item بازسازی بسته‌شده $xWW^{T}$؛
 \item حذف جریمه‌ی نرم متعامدسازی؛
 \item حفظ متعامدسازی با \QR{} سخت دوره‌ای؛
 \item استفاده از ترم‌های کوواریانسی خام؛
 \item عدم نیاز به زیان کنترل مقیاس، چون \QR{} طول ستون‌ها را کنترل می‌کند.
\end{itemize}

تفاوت اصلی این خانواده با روش ۱ استاندارد آن است که
\begin{equation}
 \lambda_O=0
\end{equation}
است و کنترل هندسی فقط به‌وسیله‌ی بازگردانی سخت اقلیدسی انجام می‌شود.

\subsection{اعضای خانواده}

\begin{center}
\begin{tabularx}{0.96\textwidth}{|C{1.5cm}|C{3.3cm}|C{4cm}|X|}
\hline
\textbf{روش} & \textbf{ترم‌های فعال} & \textbf{تابع هدف کوتاه} & \textbf{توضیح} \\
\hline
۴ & \Dterm{}+\Iterm{}+\Nterm{}+\Sterm{} & $D+I+N+S$ & نسخه‌ی کامل خانواده \\
\hline
۵ & \Dterm{}+\Iterm{}+\Sterm{} & $D+I+S$ & حذف سرکوب نویز \\
\hline
۶ & \Dterm{}+\Iterm{}+\Nterm{} & $D+I+N$ & حذف حفظ سیگنال \\
\hline
۷ & \Dterm{}+\Iterm{} & $D+I$ & هسته‌ی نویززدایی و ناوردایی \\
\hline
\end{tabularx}
\end{center}

\section{خانواده‌ی روش‌های ۸ تا ۱۱: مقداردهی روش ۲ و آموزش آزاد}

\subsection{تعریف خانواده}

در این خانواده، خروجی دقیق و متریک‌متعامد روش ۲ به‌عنوان وزن اولیه‌ی مرحله‌ی گرادیانی استفاده می‌شود:
\begin{equation}
 W_2
 \longrightarrow
 \text{آموزش آزاد روش ۱}
 \longrightarrow
 W_{\mathrm{final}}.
\end{equation}

منظور از آموزش آزاد این است که پس از شروع از $W_2$، هیچ بازگردانی اقلیدسی یا متریکی روی $W$ اعمال نمی‌شود. بنابراین مدل می‌تواند هم زیرفضا و هم مقیاس و زاویه‌ی ستون‌ها را تغییر دهد.

ویژگی‌های مشترک خانواده:
\begin{itemize}
 \item مقداردهی اولیه با خروجی دقیق \NPtwo{}؛
 \item عدم اجرای \QR{} اقلیدسی روی $W_2$؛
 \item عدم استفاده از جریمه‌ی نرم متعامدسازی؛
 \item عدم استفاده از بازگردانی سخت؛
 \item بازسازی با پروژکتور عمومی؛
 \item استفاده از ترم‌های کوواریانسی خام؛
 \item فعال‌کردن زیان کنترل مقیاس با وزن حداقل $0.1$.
\end{itemize}

تابع هدف عمومی این خانواده به‌صورت زیر است:
\begin{equation}
 \mathcal{L}
 =
 \lambda_DD
 +\lambda_II
 +\lambda_NN
 +\lambda_SS
 +\lambda_{\mathrm{scale}}\mathcal{L}_{\mathrm{scale}},
\end{equation}
که بنا بر عضو خانواده، وزن $N$ یا $S$ می‌تواند صفر باشد.

\subsection{اعضای خانواده}

\begin{center}
\begin{tabularx}{0.96\textwidth}{|C{1.5cm}|C{3.3cm}|C{4cm}|X|}
\hline
\textbf{روش} & \textbf{ترم‌های فعال} & \textbf{تابع هدف کوتاه} & \textbf{توضیح} \\
\hline
۸ & \Dterm{}+\Iterm{}+\Nterm{}+\Sterm{} & $D+I+N+S$ & آموزش آزاد کامل از $W_2$ \\
\hline
۹ & \Dterm{}+\Iterm{}+\Sterm{} & $D+I+S$ & بدون سرکوب نویز \\
\hline
۱۰ & \Dterm{}+\Iterm{}+\Nterm{} & $D+I+N$ & بدون حفظ سیگنال \\
\hline
۱۱ & \Dterm{}+\Iterm{} & $D+I$ & فقط هسته‌ی گرادیانی روی مقداردهی روش ۲ \\
\hline
\end{tabularx}
\end{center}

\subsection{معنای هندسی خانواده}

این خانواده تلاش نمی‌کند پاسخ روش ۲ را ثابت نگه دارد. $W_2$ فقط یک نقطه‌ی شروع نویزآگاه است. بازسازی عمومی باعث می‌شود حتی اگر ستون‌ها غیرمتعامد باشند، بازسازی همچنان بر گستره‌ی ستونی جاری انجام شود. زیان کنترل مقیاس نیز از تغییرات نامحدود طول ستون‌ها جلوگیری می‌کند، بدون آنکه آن‌ها را مجبور به طول واحد کند.

\section{خانواده‌ی روش‌های ۱۲ تا ۱۵: مقداردهی روش ۲ و بازگردانی متریکی}

\subsection{تعریف خانواده}

این خانواده نیز از خروجی دقیق روش ۲ شروع می‌کند، اما هندسه‌ی نویزآگاه آن را در طول آموزش حفظ می‌کند:
\begin{equation}
 W_2
 \longrightarrow
 \text{آموزش گرادیانی}
 \longrightarrow
 \text{بازگردانی در متریک }B.
\end{equation}

ماتریس متریک از کوواریانس نویز همان نماهای آموزشی ساخته می‌شود:
\begin{equation}
 B=\Sigma_n+\rho I.
\end{equation}

پس از هر بازه‌ی مشخص از گام‌های بهینه‌سازی، بازگردانی زیر اجرا می‌شود:
\begin{equation}
 W\leftarrow W\left(W^{T}BW\right)^{-1/2}.
\end{equation}

ویژگی‌های مشترک خانواده:
\begin{itemize}
 \item مقداردهی اولیه با خروجی دقیق \NPtwo{}؛
 \item بازسازی با پروژکتور عمومی؛
 \item عدم استفاده از \QR{} اقلیدسی؛
 \item عدم استفاده از جریمه‌ی نرم متعامدسازی؛
 \item حفظ سخت شرط $W^{T}BW=I$ با بازگردانی متریکی؛
 \item استفاده از ترم‌های کوواریانسی خام؛
 \item فعال‌کردن زیان کنترل مقیاس با وزن حداقل $0.05$.
\end{itemize}

در این خانواده، متریک $B$ یک‌بار از کل نماهای آموزشی محاسبه و برای بازگردانی ثابت نگه داشته می‌شود. کوواریانس مورد استفاده در ترم سرکوب نویز، مطابق سازوکار روش ۱، در حالت پیش‌فرض با \EMA{} دسته‌ها به‌روزرسانی می‌شود.

\subsection{اعضای خانواده}

\begin{center}
\begin{tabularx}{0.96\textwidth}{|C{1.5cm}|C{3.3cm}|C{4cm}|X|}
\hline
\textbf{روش} & \textbf{ترم‌های فعال} & \textbf{تابع هدف کوتاه} & \textbf{توضیح} \\
\hline
۱۲ & \Dterm{}+\Iterm{}+\Nterm{}+\Sterm{} & $D+I+N+S$ & آموزش کامل با حفظ هندسه‌ی متریکی \\
\hline
۱۳ & \Dterm{}+\Iterm{}+\Sterm{} & $D+I+S$ & بدون سرکوب نویز \\
\hline
۱۴ & \Dterm{}+\Iterm{}+\Nterm{} & $D+I+N$ & بدون حفظ سیگنال \\
\hline
۱۵ & \Dterm{}+\Iterm{} & $D+I$ & هسته‌ی گرادیانی با بازگردانی متریکی \\
\hline
\end{tabularx}
\end{center}

\subsection{تفاوت با خانواده‌ی آزاد}

در خانواده‌ی آزاد، $W^{T}BW$ می‌تواند در طول آموزش از ماتریس همانی فاصله بگیرد. در این خانواده، همان آزادی فقط در فضای مماس میان دو بازگردانی وجود دارد و سپس شرط متریکی دوباره برقرار می‌شود. بنابراین، روش ۲ علاوه بر مقداردهی اولیه، تعریف هندسه‌ی مجاز آموزش را نیز تعیین می‌کند.

\section{خانواده‌ی روش‌های ۱۶ تا ۱۹: \QR{} خروجی روش ۲ و سپس روش ۱}

\subsection{تعریف خانواده}

در این خانواده، خروجی روش ۲ ابتدا به یک پایه‌ی اقلیدسی‌متعامد تبدیل می‌شود:
\begin{equation}
 W_2=QR,
 \qquad
 W_0=Q.
\end{equation}

سپس روش ۱ با مقداردهی $Q$ و \QR{} سخت دوره‌ای آموزش داده می‌شود:
\begin{equation}
 W_2
 \xrightarrow{\mathrm{Euclidean\;QR}}
 Q
 \longrightarrow
 \text{آموزش روش ۱ با هندسه‌ی اقلیدسی}.
\end{equation}

ویژگی‌های مشترک خانواده:
\begin{itemize}
 \item استخراج زیرفضای اولیه با \NPtwo{}؛
 \item تبدیل خروجی روش ۲ به پایه‌ی اقلیدسی با \QR{}؛
 \item حفظ گستره‌ی ستونی اولیه و حذف مقیاس و هندسه‌ی متریکی داخل آن؛
 \item بازسازی بسته‌شده $xWW^{T}$؛
 \item حذف جریمه‌ی نرم متعامدسازی؛
 \item حفظ سخت شرط $W^{T}W=I$ با \QR{} دوره‌ای؛
 \item عدم نیاز به زیان کنترل مقیاس.
\end{itemize}

اگر $W_2$ تمام‌رتبه باشد، عملیات \QR{} گستره‌ی ستونی آن را حفظ می‌کند:
\begin{equation}
 \operatorname{span}(W_2)=\operatorname{span}(Q),
\end{equation}
اما مختصات و مقیاس‌های خاص حل تعمیم‌یافته را کنار می‌گذارد.

\subsection{اعضای خانواده}

\begin{center}
\begin{tabularx}{0.96\textwidth}{|C{1.5cm}|C{3.3cm}|C{4cm}|X|}
\hline
\textbf{روش} & \textbf{ترم‌های فعال} & \textbf{تابع هدف کوتاه} & \textbf{توضیح} \\
\hline
۱۶ & \Dterm{}+\Iterm{}+\Nterm{}+\Sterm{} & $D+I+N+S$ & نسخه‌ی کامل روی پایه‌ی \QR{}شده‌ی روش ۲ \\
\hline
۱۷ & \Dterm{}+\Iterm{}+\Sterm{} & $D+I+S$ & بدون سرکوب نویز \\
\hline
۱۸ & \Dterm{}+\Iterm{}+\Nterm{} & $D+I+N$ & بدون حفظ سیگنال \\
\hline
۱۹ & \Dterm{}+\Iterm{} & $D+I$ & هسته‌ی روش ۱ روی زیرفضای اقلیدسی روش ۲ \\
\hline
\end{tabularx}
\end{center}

\section{مقایسه‌ی ساختاری خانواده‌های فعلی}

\begin{longtable}{|C{1.8cm}|C{3.1cm}|C{3cm}|C{3.2cm}|C{3.3cm}|}
\hline
\textbf{روش‌ها} & \textbf{مقداردهی اولیه} & \textbf{قید هندسی سخت} & \textbf{بازسازی} & \textbf{کنترل مقیاس} \\
\hline
\endfirsthead
\hline
\textbf{روش‌ها} & \textbf{مقداردهی اولیه} & \textbf{قید هندسی سخت} & \textbf{بازسازی} & \textbf{کنترل مقیاس} \\
\hline
\endhead
۴--۷ & تصادفی متعامد & $W^{T}W=I$ با \QR{} & $xWW^{T}$ & ناشی از \QR{} \\
\hline
۸--۱۱ & خروجی دقیق روش ۲ & بدون قید سخت & پروژکتور عمومی & زیان مرتبه‌ی چهار، وزن حداقل $0.1$ \\
\hline
۱۲--۱۵ & خروجی دقیق روش ۲ & $W^{T}BW=I$ & پروژکتور عمومی & زیان مرتبه‌ی چهار، وزن حداقل $0.05$ \\
\hline
۱۶--۱۹ & \QR{} خروجی روش ۲ & $W^{T}W=I$ با \QR{} & $xWW^{T}$ & ناشی از \QR{} \\
\hline
\end{longtable}

\subsection{چیزی که هر خانواده از روش ۲ حفظ می‌کند}

\begin{itemize}
 \item \textbf{روش‌های ۴ تا ۷:} هیچ بخش از روش ۲ را استفاده نمی‌کنند و یک مرجع مستقل برای روش ۱ هستند.
 \item \textbf{روش‌های ۸ تا ۱۱:} هم زیرفضای اولیه و هم مقیاس و هندسه‌ی داخلی $W_2$ را در نقطه‌ی شروع حفظ می‌کنند، ولی اجازه می‌دهند همه‌ی آن‌ها در آموزش تغییر کنند.
 \item \textbf{روش‌های ۱۲ تا ۱۵:} خروجی اولیه و هندسه‌ی متریکی روش ۲ را حفظ می‌کنند و فقط تغییراتی را می‌پذیرند که پس از بازگردانی با شرط $W^{T}BW=I$ سازگار باشند.
 \item \textbf{روش‌های ۱۶ تا ۱۹:} فقط گستره‌ی ستونی روش ۲ را در نقطه‌ی شروع حفظ می‌کنند و مقیاس و هندسه‌ی متریکی آن را با \QR{} اقلیدسی حذف می‌کنند.
\end{itemize}

\section{فهرست نهایی روش‌های توضیح‌داده‌شده}

\begin{center}
\begin{tabularx}{0.98\textwidth}{|C{1.4cm}|X|C{3.5cm}|}
\hline
\textbf{شماره} & \textbf{تعریف} & \textbf{ترم‌های فعال} \\
\hline
۱ & \NPone{} استاندارد با بازسازی بسته‌شده، جریمه‌ی نرم اقلیدسی و \QR{} سخت & $D+I+N+S+O$ \\
\hline
۲ & \NPtwo{}؛ حل بسته‌ی مقدارویژه‌ی تعمیم‌یافته با متعامدسازی در متریک نویز & حل بسته \\
\hline
۴ & روش ۱ با \QR{} سخت و بدون جریمه‌ی نرم & $D+I+N+S$ \\
\hline
۵ & خانواده‌ی روش ۱ با \QR{}، بدون سرکوب نویز & $D+I+S$ \\
\hline
۶ & خانواده‌ی روش ۱ با \QR{}، بدون حفظ سیگنال & $D+I+N$ \\
\hline
۷ & خانواده‌ی روش ۱ با \QR{}، هسته‌ی دوترمی & $D+I$ \\
\hline
۸ & خروجی دقیق روش ۲ و آموزش آزاد کامل & $D+I+N+S$ \\
\hline
۹ & خروجی دقیق روش ۲ و آموزش آزاد، بدون سرکوب نویز & $D+I+S$ \\
\hline
۱۰ & خروجی دقیق روش ۲ و آموزش آزاد، بدون حفظ سیگنال & $D+I+N$ \\
\hline
۱۱ & خروجی دقیق روش ۲ و آموزش آزاد، هسته‌ی دوترمی & $D+I$ \\
\hline
۱۲ & خروجی دقیق روش ۲ و آموزش کامل با بازگردانی متریکی & $D+I+N+S$ \\
\hline
۱۳ & خانواده‌ی متریکی، بدون سرکوب نویز & $D+I+S$ \\
\hline
۱۴ & خانواده‌ی متریکی، بدون حفظ سیگنال & $D+I+N$ \\
\hline
۱۵ & خانواده‌ی متریکی، هسته‌ی دوترمی & $D+I$ \\
\hline
۱۶ & \QR{} خروجی روش ۲ و سپس روش ۱ کامل & $D+I+N+S$ \\
\hline
۱۷ & خانواده‌ی \QR{} روش ۲، بدون سرکوب نویز & $D+I+S$ \\
\hline
۱۸ & خانواده‌ی \QR{} روش ۲، بدون حفظ سیگنال & $D+I+N$ \\
\hline
۱۹ & خانواده‌ی \QR{} روش ۲، هسته‌ی دوترمی & $D+I$ \\
\hline
\end{tabularx}
\end{center}

\section{جمع‌بندی تعریف‌ها}

روش ۱ یک نگاشت خطی چندهدفه است که با نویززدایی، ناوردایی تقابلی، سرکوب کوواریانس نویز، حفظ واریانس سیگنال و کنترل اقلیدسی هندسه آموزش می‌بیند. روش ۲ همان مسئله را از دید آمار مرتبه‌ی دوم حل می‌کند و زیرفضایی را با بیشینه‌کردن نسبت سیگنال به نویز و قید متعامدسازی در متریک نویز به‌دست می‌آورد.

خانواده‌های فعلی چهار نوع استفاده از این دو ایده را پوشش می‌دهند: آموزش مستقل روش ۱ در هندسه‌ی اقلیدسی، شروع از پاسخ روش ۲ و آموزش آزاد، شروع از پاسخ روش ۲ با حفظ هندسه‌ی متریکی، و حفظ فقط زیرفضای روش ۲ پس از \QR{} اقلیدسی. در هر خانواده، چهار حالت ترمی یکسان تعریف شده‌اند تا نقش دو ترم سرکوب نویز و حفظ سیگنال مستقل از انتخاب هندسه بررسی شود.

\end{document}
